% =====================================================================
% FINESST-25 — FILE 2: "EXPERTISE AND RESOURCES — NOT ANONYMIZED"
% Full draft. Real names REQUIRED here.
% RRS <= 1 page. Acknowledgement <= 150 words.
% After compiling: append the four PDFs exported from the NASA docx
% templates (biosketches + C&P for PI and FI) — see drafts in
% ../accompanying/.
% =====================================================================
\documentclass[12pt,letterpaper]{article}
\usepackage[margin=1in]{geometry}
\usepackage{newtxtext,newtxmath}
\usepackage{booktabs}
\usepackage{xcolor}
\newcommand{\todo}[1]{\textcolor{red}{[TODO: #1]}}

\begin{document}

\begin{center}{\large\bfseries Expertise and Resources --- Not Anonymized}\\[3pt]
{\itshape Mapping Reionization-Era Ionized Bubbles by Recovering
Lyman-$\alpha$ Information from Low-Resolution JWST and Roman
Spectroscopy}
\end{center}

% ---- (a) team list — must be page 1 ----
\section*{1. Team Members, Affiliations, and Roles}
\begin{center}
\begin{tabular}{lll}
\toprule
Name & Affiliation & Role \\
\midrule
Prof.\ Bahram Mobasher & University of California, Riverside & PI \\
Aryana Haghjoo & University of California, Riverside & FI (PhD student) \\
Dr.\ Shoubaneh Hemmati & IPAC, California Institute of Technology & Collaborator (Mentor 1) \\
\bottomrule
\end{tabular}
\end{center}

\section*{2. Team Expertise and Responsibilities}

\textbf{Aryana Haghjoo (FI)} developed the physics-informed spectral
super-resolution framework at the core of this proposal and is first
author of the paper presenting it (Haghjoo et al.\ 2026,
arXiv:2603.18357, submitted to ApJ). The FI also built the end-to-end
Roman grism simulation pipeline (extraction, contamination modeling,
training-set generation from the Wang et al.\ 2022 IRSA products) and the
Lyman-$\alpha$ mock-spectrum pipeline used in Task 2. The FI's M.Sc.\
research (McGill University) on parameter estimation of the global 21\,cm
signal provides direct background in Epoch-of-Reionization physics and
Bayesian inference. The FI will carry out all proposed research and lead
all publications.

\textbf{Prof.\ Bahram Mobasher (PI)} is a Distinguished Professor at the
University of California, Riverside, with over three decades of
experience in galaxy evolution, deep multiwavelength surveys,
photometric-redshift and SED-fitting methodology, and space-based survey
science (HST, Spitzer, JWST, Euclid, Roman preparation). The PI
supervises the FI's PhD thesis, directs the overall research program,
and mentors the FI in survey science and scientific writing.

\textbf{Dr.\ Shoubaneh Hemmati (Collaborator, Mentor 1)} is a Research
Scientist at IPAC/Caltech specializing in machine-learning methods for
galaxy surveys and spectroscopy; co-developer of the super-resolution
framework (second author of Haghjoo et al.\ 2026). Dr.\ Hemmati advises
on model architecture, validation methodology, and Roman survey
spectroscopy, and co-mentors the FI (weekly meetings, plus in-person
working sessions at IPAC approximately every two weeks).

\section*{3. Summary of Work Effort}
\begin{center}
\begin{tabular}{lll}
\toprule
Name (role) & Contribution & Effort \\
\midrule
Aryana Haghjoo (FI) & All analysis, pipelines, publications & 1.0 FTE \\
Bahram Mobasher (PI) & Thesis supervision, survey science guidance & no cost requested \\
Shoubaneh Hemmati (Mentor 1) & ML methodology co-mentoring & no cost requested \\
\bottomrule
\end{tabular}
\end{center}

\section*{4. Facilities, Equipment, and Other Resources}

\textbf{Computing.} The PI's research group at UCR operates a dedicated
GPU system equipped with an NVIDIA GeForce RTX 5090 (32~GB), supplemented
by the University of California, Riverside High Performance Computing
Center (HPCC), which provides shared GPU nodes and multi-TB project
storage. Model training for this project ($\sim$2.5M-parameter networks,
datasets $\lesssim$10 GB) is well within these resources; no NASA HEC
allocation is required.

\textbf{Data access.} All observational inputs are public: JWST/NIRSpec
JADES data releases (MAST) and Roman grism simulation products (Wang et
al.\ 2022; IRSA). No letters of resource support are required for these
public archives.

\textbf{Simulations.} Radiative-transfer post-processed hydrodynamic
simulations of patchy reionization (Lyman-$\alpha$ transmission
sightlines with matched ionization cubes) are provided through an
ongoing collaboration with Prof.\ Anson D'Aloisio (UC Riverside); the
FI already holds the simulation data and has built the training-label
pipeline on them, so continued access is assured and no letter of
resource support is required. Public simulation suites (THESAN,
SPHINX$_{20}$) provide independent training sets.

% ---- RRS — FI-authored, MAX 1 PAGE ----
\clearpage
\section*{5. Research Readiness Statement (Future Investigator)}

I am a PhD student in Astronomy at the University of California,
Riverside (enrolled since September 2024 in the M.Sc./Ph.D.\ track;
Ph.D.\ program since September 2025), with expected graduation in June
2030 --- fully spanning the proposed three-year award. My path to this
project combines the three ingredients it requires: reionization physics,
machine-learning methodology, and space-based spectroscopy.

\emph{Reionization and statistical inference.} My M.Sc.\ in Physics at
McGill University (2021--2023) was devoted to the Epoch of Reionization:
I built an MCMC parameter-estimation pipeline for the global 21\,cm
signal using ARES simulations and EDGES data, giving me a working
foundation in EoR astrophysics and Bayesian inference --- the physics and
statistics underlying Tasks 1 and 2.

\emph{Machine learning for spectroscopy.} I am first author of
``Learning to See Sharper: A Physics-Informed Artificial Intelligence
Framework for Super-Resolving Galaxy Spectra'' (arXiv:2603.18357,
submitted to ApJ), the framework at the center of this proposal: I
designed the architecture, training strategy, and validation
methodology, trained it on 1,187 JADES prism--grating pairs, and released
the code and checkpoints publicly. I have since ported the framework to
simulated Roman grism data, building the full extraction and
training pipeline (17,000+ spectral pairs to date) and the
prior-dominance audit described in the S/T/M. My ML background began with
undergraduate research on variable-star classification and coursework in
machine learning for physics; my current graduate coursework includes
galaxy formation, advanced cosmology, radiative processes, and
observational astrophysics.

\emph{Observational experience.} I work with JWST/NIRSpec data daily
(JADES DR3/DR4); I am a Co-I of an accepted JWST Cycle 5 GO program
(NIRSpec/IFU observations of extremely luminous $z\sim9$ galaxies,
unfunded role); and I have contributed to spectroscopic catalog
construction from Keck data for the H2O survey. I have presented this research program at ten venues
since mid-2025, including invited talks at the CfA AstroAI seminar and
the Euclid US science meeting.

I am proficient in Python/PyTorch, high-performance computing (SLURM),
and collaborative software development, and have prototyped every
pipeline this proposal requires: the mock Lyman-$\alpha$ generation
(Task 1), the bubble-size training-label extraction from
radiative-transfer cubes (Task 2), and the Roman injection machinery
(Task 3). The proposed research is the natural integration of these
components into a single measurement, and I am prepared to execute it.

% ---- Acknowledgement — MAX 150 WORDS ----
\clearpage
\section*{6. Acknowledgement of Authorship}

I, Aryana Haghjoo, affirm that I am the primary author of this proposal.
I defined the scientific objectives, designed the technical approach,
and wrote the Science/Technical/Management section, the Open Science and
Data Management Plan, and this Expertise and Resources document. The
mentoring plan was prepared jointly with the PI, Prof.\ Bahram Mobasher,
and collaborator Dr.\ Shoubaneh Hemmati, who also provided scientific
feedback on the research plan. The proposed research builds on my
first-author work developing the spectral super-resolution framework and
represents my own initiative in extending it to reionization science.
Portions of the text, \LaTeX, and figure-production code were drafted
and revised with the assistance of a generative AI tool (Claude,
Anthropic, June--July 2026) under my direction; all scientific ideas,
methods, and conclusions are my own, and I take full responsibility for
the accuracy of the entire proposal.

% ---- Budget ----
\section*{7. Budget and Budget Narrative}

% Structure agreed with the PI (2026-07-07): salary first (GSR-1 summer
% 2027, GSR-3 after advancement to candidacy), remainder to tuition;
% travel $1,500/yr; no publication charges in Year 1; UCR treats FINESST
% FIs as fellowship holders (no F&A). Cap: $50,000/year TOTAL.

\subsection*{Period of performance}
Three years, 1 June 2027 -- 31 May 2030, beginning with the Summer 2027
quarter. (The solicitation requires a start date no later than 14 July
2027; June 1, 2027 satisfies this and covers the full summer quarter.)

\subsection*{Year 1 (June 2027 -- May 2028)}
\begin{center}
\begin{tabular}{lr}
\toprule
Item & Amount \\
\midrule
FI stipend --- Summer 2027 quarter: UC GSR Step 1, 50\% appointment
  (3 mo $\times$ \$2,995.63) & \$8,987 \\
FI stipend --- Fall--Spring: UC GSR Step 3 upon advancement to
  candidacy, 50\% (9 mo $\times$ \$3,477.96) & \$31,302 \\
Domestic travel (AAS or one topical workshop) & \$1,500 \\
Tuition and fee contribution & \$8,211 \\
\midrule
Total Year 1 (fellowship classification; no F\&A) & \$50,000 \\
\bottomrule
\end{tabular}
\end{center}

\subsection*{Years 2 and 3 (each)}
\begin{center}
\begin{tabular}{lr}
\toprule
Item & Amount \\
\midrule
FI stipend --- 12 months at UC GSR Step 3, 50\%
  (12 mo $\times$ \$3,477.96) & \$41,736 \\
Domestic travel & \$1,500 \\
Publication charges (open access, 1 paper/yr) & \$2,000 \\
Tuition and fee contribution & \$4,764 \\
\midrule
Total per year & \$50,000 \\
Three-year total & \$150,000 \\
\bottomrule
\end{tabular}
\end{center}

\subsection*{Narrative}
The requested funds support the FI only; no salary, travel, or other
costs are requested for the PI or collaborators, per the solicitation.
The University of California, Riverside administers FINESST awardees as
fellowship holders; accordingly no facilities and administrative costs
are assessed and the full annual amount supports the FI directly.

The budget prioritizes the FI's stipend so that the FI is released from
teaching duties for the entire award period. Stipend amounts follow the
published University of California Graduate Student Researcher salary
scale (UCOP Table 22, rates effective 10/1/2025: Step 1
\$5,991.25/month, Step 3 \$6,955.92/month, full-time equivalent) at a
50\% appointment (20 hours per week) --- the maximum employment level
for the FI as an international student and the standard UC graduate
appointment: Step 1 during the Summer 2027 quarter and Step 3 from Fall
2027 onward, reflecting the FI's advancement to Ph.D.\ candidacy. Should the UC salary scale be adjusted
during the award period, the stipend will follow the prevailing scale
and the tuition contribution line will absorb the difference within the
annual cap.

The remaining annual amount contributes toward tuition and fees. Per the
University of California 2025--26 fee schedule, graduate academic
tuition and the student services fee total \$4,810 per quarter (three
quarters per year; the summer quarter carries no tuition), plus
campus-based fees and health insurance; nonresident supplemental tuition
(\$5,034 per quarter) is assessed only until advancement to candidacy,
after which it is \$0 for the duration of the award. The balance of
tuition and fees not covered by this award is paid from the PI's
research funds --- a committed cost share that ensures the project is
fully resourced within the FINESST annual cap.

Travel (\$1,500 per year) supports presentation of results at one
national meeting (e.g., AAS) or one topical workshop per year. The FI's
regular in-person working sessions with Dr.\ Hemmati at IPAC involve
only local travel and require no budgeted funds. Publication charges
(\$2,000 in Years 2 and 3 only) cover one open-access article per year,
matching the publication schedule in the S/T/M: the Year-1 methods paper
is expected to reach publication charges in Year 2, followed by the
flagship measurement and Roman forecast papers.

% =====================================================================
% APPEND (exported to PDF from ../accompanying/ drafts + NASA docx):
%   biosketch_PI.pdf, biosketch_FI.pdf, [optional: Hemmati as mentor],
%   CP_PI.pdf, CP_FI.pdf
% NO transcripts, NO recommendation letters.
% =====================================================================

\end{document}
